%% Narrow-screen preview build.
%% Same manuscript, re-flowed onto a ~110 mm wide single-column page so it
%% stays legible when a PDF viewer shows it fit-to-width in a pane that is
%% narrower than a quarter of the screen (side-by-side writing/preview).
%%
%% Build:  ./build-narrow.sh     (or: latexmk -pdf 00_Article_Narrow.tex)
%%
%% The other two layouts are unchanged:
%%   \documentclass[twocolumn]{theme/bioRxiv}   % bioRxiv preprint look
%%   \documentclass[submit]{theme/bioRxiv}      % journal submission look
\documentclass[submit,narrow]{theme/bioRxiv}

%% forces sans serif
\renewcommand\familydefault{\sfdefault}

%% for this example - not needed for manuscript
\usepackage{lipsum}

%% setup specific to this preprint -- add or remove as you need.
\hyphenation{empire} % prevents word "empire" being split across lines
\sisetup{range-units = single,
        separate-uncertainty = true,
        multi-part-units = single,
        range-phrase = --} % units setup
\usepackage{multirow} % allows merging of cells in tables

%% `\narrowfit{...}` (defined by the class) scales a too-wide table down to the
%% column; it is a no-op in the other layouts.

%% reading aids for a narrow pane: real PDF outline + clickable links
\hypersetup{
  pdftitle={Description of an amazing research project},
  pdfauthor={First Scientist; Second Doctor; Third Professor},
  bookmarksnumbered=true,
  bookmarksopen=true,
  bookmarksopenlevel=1,
  pdfstartview={FitH},
}

%% the document starts here

\begin{document}

%% comment out any files that you don't want to include as sections
 \leadauthor{Scientist}

\title{Description of an amazing research project}
\shorttitle{Running title here}

\author[1]{First Scientist \orcidlink{0000-0001-0000-0000}}
\author[2]{Second Doctor \orcidlink {000-0002-0000-0000}}
\author[1,\Letter]{Third Professor \orcidlink {000-0003-0000-0000}}
\affil[1]{A University, Academic Street, Learningtown, UK}
\affil[2]{B Institute, Chalk Road, Blackboardville, USA}
\date{}

\maketitle

\begin{abstract}
Abstract of the paper goes here.
\lipsum[1]
\end{abstract}

\begin{keywords}
keyword1 | keyword2 | keyword3
\end{keywords}

\begin{corrauthor}
third.professor\at awesome.ac.uk
\end{corrauthor}

\section*{Introduction}\label{s:introduction}

\lipsum{2-6}

\section*{Results}\label{s:results}

\subsection*{Citations and full size figures with legends underneath}

Text is added like this
This is a reference to a published paper \citep{watson_molecular_1953}.
We can cite other things too \citep{tipton_complexities_2019,zheng_genome_2011,alberts_molecular_2002}

This is a new paragraph.
New sentences on a new line.
New sentences on a new line.

% this is how to add a comment
This is a new result.
% this is how to add a figure with the name cells.
As you can see (Figure \ref{fig:cells}).

% full size figure is figure*
\begin{figure*}
\centering
\includegraphics[width=0.75\linewidth]{figures/temp.png}
\caption{\textbf{These are cells.}\\
(\textbf{A}) This is a regular figure with a legend as a caption underneath. Inset: 3X zoom. Scale bar, \SI{10}{\micro\meter}.}
\label{fig:cells}
\end{figure*}

It is possible to add a one-column Figure like this (Figure \ref{fig:nucleus}).
To add Supplementary Figures you can do either of these things and have them at the end of the end of the paper (Supplementary Figure \ref{suppfig:endosome}).
Or like this (Supplementary Figure \ref{videosupp:lysosome}).

\subsection*{Working with tables and equations}

References to a table can be included here like this (Table \ref{tab:t1}).
If you have a Supplementary Table sections you can reference one of those tables like this (Supplementary Table \ref{supptab:st1})

\begin{table}[h] % h means "here"
    \centering
    \begin{tabular}{l|l}
         Base & Letter\\
         \hline
         Adenine & A\\
         Cytosine & C\\
         Guanine & G\\
         Thymidine & T\\
    \end{tabular}
    \caption{The DNA alphabet}
    \label{tab:t1}
\end{table}

To include an equation, you can write in-line math like this, $e = mc^2$, or you can reference an equation like this (Equation \ref{eqn:euler}).

\begin{equation}
    \centering
    e^{\pi i} + 1 = 0
    \label{eqn:euler}
\end{equation}

\lipsum[10]

\subsection*{Subsections are written like this}

\lipsum[11]

% one-column size figure is figure
% note that when using `submit` class option, the width will refer to the full page, so adjust accordingly
\begin{figure}
\centering
\includegraphics[width=0.75\linewidth]{figures/temp.png}
\caption{\textbf{This is a nucleus.}\\
(\textbf{A}) This is a one-column figure with a legend as a caption underneath.}
\label{fig:nucleus}
\end{figure}

\lipsum[12]

\subsection*{Another subsection}

\lipsum[13-14]

\subsection*{Another subsection}

\lipsum[13-14]

\subsection*{Another subsection}

\lipsum[13-14]

\section*{Discussion}\label{s:discussion}

This is the discussion section where you wax lyrical about your findings.
You can put your work in the context of other published work \citep{brenner_uga:_1967}.

\lipsum[100-104]

\section*{Methods}\label{s:methods}

\subsection*{Molecular biology}

Details of plasmids are usually first.
Followed by cell biology section.
We have special units defined for: molar, units, spin speed, electron and Angstrom, i.e. \SI{1}{\Molar} sucrose, which can also be written \SI{1}{\molar} sucrose, \SI{10}{\Units\per\milli\litre} restriction enzyme, centrifuge at \SI{10000}{\gee} or \SI{800}{\rpm}, this at \SI{100}{\kilo\electron}, a resolution of \SI{1.8}{\angstrom}.
Otherwise use \texttt{siunitx} for everything else, \SI{10}{\micro\metre} and \SI{37}{\degreeCelsius} and what-not.
There are some gotchas with using \texttt{siunitx} for units, so check the documentation.
For example, the unit for minute is \texttt{\textbackslash minute} and not \texttt{\textbackslash min}, which is a reserved word in LaTeX.
If you are seeing an error in compilation, check the units you are using.

You can use ``chemical entities'' (ce) for chemical formulas, such as \ce{Na2HPO4} and \ce{H2O}.

\subsection*{Cell biology}

The general rule is that the methods subsections should be in the order that the results are presented, and in the order that the experiments were performed.

\lipsum[80]

\begin{acknowledgements}
We thank Prof Ersatz for useful discussion and all members of our research groups for constructive feedback.
The work was supported by grants from this charity and that foundation (Award \#666).

For the purpose of open access, the authors have applied a Creative Commons Attribution (CC BY) license to any Author Accepted Manuscript version arising
\end{acknowledgements}

\begin{contributions}
FS, all experimental work, data analysis, programming, first draft;
SD, funding acquisition and supervision;
TP, data analysis, programming, wrote the manuscript.
\end{contributions}

\begin{interests}
The authors declare no conflict of interest.
\end{interests}

\section*{References}
\bibliographystyle{theme/bxv_abbrvnat}
\bibliography{refs.bib}
\include{02_Article_Supplementary}
\include{03_Article_SupplementaryVideos}
%  CONFIGURE NEW SINGLE-PAGE FORMAT 

\onecolumn % go back to one column
\fancyhead{} % make sure we get no headers
\renewcommand{\floatpagefraction}{0.1}
\lfoot[\bSupInf]{\dAuthor}
\rfoot[\dAuthor]{\cSupInf}
\newpage

\captionsetup*{format=largeformat} % make legend slightly larger than in the paper
\setcounter{table}{0} % reset table counter for Supplementary Video Figures
%\setcounter{page}{1} % reset page count
\makeatletter 
\renewcommand{\thetable}{S\@arabic\c@table} % make Table legend start with Table S
\makeatother

%  MAIN TEXT 
\newpage
\section*{Supplementary Tables}

\bigskip

\begin{table}[!h]
    \centering
    \small
    %% \narrowfit scales the table to the line width on the narrow-screen page
    %% and is a no-op in the preprint/submission layouts
    \narrowfit{%
    \begin{tabular}{p{3cm}|p{2.2cm} p{2.2cm} p{2.2cm} p{2.2cm} p{2.2cm}}
        Movie & US Release Date & Director & Screenwriter(s) & Story by & Producer\\
        \hline
        Episode IV – A New Hope & May 25, 1977 & \multicolumn{3}{c}{George Lucas} & Gary Kurtz\\
        Episode V – The Empire Strikes Back & May 21, 1980 & Irvin Kershner & Leigh Brackett and Lawrence Kasdan & George Lucas & Gary Kurtz\\
        Episode VI – Return of the Jedi & May 25, 1983 & Richard Marquand & Lawrence Kasdan and George Lucas & George Lucas & Howard Kazanjian\\
    \end{tabular}}
    \caption{\textbf{Original ``Star Wars'' Trilogy}\\
    Extracted from \url{https://en.wikipedia.org/wiki/List_of_Star_Wars_films}.}
    \label{supptab:st1}
\end{table}





\end{document}
